\section{Specification and Roadmap Summary}
\label{sec:spec}
The companion \emph{Technical Specification} and \emph{Development Roadmap}
are standalone documents built from the same sources as this report; this
section summarises them so that the report is self-contained.

\subsection{Specification}
The specification uses RFC~2119 keywords and gives every requirement an
identifier: functional (FR), non-functional (NFR), interface (IF), model
(MR), training (TR), security (SR), operations (OR) and evaluation (ER). The
requirement catalogue has \NREQ\ entries; each is traced to the result or
design decision that motivates it, to the component that implements it and
to the test that verifies it. Key normative content:
\begin{itemize}
\item \textbf{Gate (FR-010--FR-014).} The controller MUST compute the
decision of Theorem~\ref{thm:gate} from a calibrated estimator and a
configured cost ratio, MUST never keep a verifiable draft that failed
verification, and MUST write a decision record with the features, estimates
and thresholds used.
\item \textbf{Deliberation (FR-020--FR-024).} Sequential verified sampling
with configured caps $N$, $R$ and budget $b$ per domain derived from the
Phase-0 fit; draft-blind by default; budget forcing on exhaustion; abstention
with notice.
\item \textbf{Verification (FR-030--FR-033).} Exact verifiers accept at
threshold $1.0$; every verifier publishes $(r,\varphi)$; assurance levels are
attached to every reply; consensus never writes training data.
\item \textbf{Self-model (FR-040--FR-042).} A $5\%$ audit stream; unmeasured
capabilities reported as unmeasured; statistics keyed by model and adapter
hash.
\item \textbf{Distillation (TR-001--TR-008).} Data schema, filtering, Unsloth
LoRA configuration, GPU path, promotion gate with confidence bounds and
held-out verifier, human sign-off, GGUF export.
\item \textbf{Targets (NFR-001--NFR-012).} Hard latency bound from
Lemma~\ref{lem:termination}; fast-path first token under two seconds and
150-token reply within twelve seconds at p95 on the reference machine;
resident memory under 8\,GB for Profile~A; reproducible builds; structured
telemetry. All targets are labelled as targets.
\item \textbf{Interfaces (IF-001--IF-010).} OpenAPI 3.1 for the gateway
with SSE streaming and a \texttt{reasoning\_mode} field on every reply; JSON
Schemas for the request, reply, decision record and trace; the backend
contract to vLLM and llama-server with retry, timeout and backpressure
policy.
\end{itemize}
Machine-readable artefacts accompany the document: \texttt{openapi.yaml},
JSON Schemas, a sample vLLM launch configuration and llama-server command,
a Cargo workspace skeleton and a sample configuration file, all validated
with standard tools.

\subsection{Roadmap}
Fig.~\ref{fig:gantt} shows the seven phases and their gates; effort and dates
are estimates. Phase~0 (foundations, weeks 0--3) delivers the serving
baseline on the reference machine and the $4\,600$-task measurement run.
Phase~1 (CPU MVP, weeks 3--8) delivers the gateway, both paths behind a
fixed-rule gate, the chat UI and the minimal verifier hub, which is the
smallest build that demonstrates the idea end to end. Phase~2 (gate and
loop, weeks 8--13) replaces the fixed rule with the calibrated gate and adds
the verified retry loop and the audit stream; its exit gate is the
non-inferiority test at three points. Phase~3 (fine-tuning, weeks 12--18)
runs the first distillation cycle on a rented GPU under the promotion gate.
Phase~4 (evaluation suite and performance console, weeks 15--21) completes
the dashboard and the load tests. Phase~5 (hardening and optional GPU
scale-out, weeks 20--25) adds security hardening, the vLLM CUDA profile and
the Qwen3.5 upgrade decision. Phase~6 (release and publication, weeks
25--28) delivers the one-point non-inferiority gate, documentation and the
research write-up. Phase~7 (CPU inference performance, added after the
first six phases were implemented) attributes latency to stages, tunes the
engine under interleaved A/B tests, adds early escalation, and prepares
EAGLE-3 speculative heads and repacked weight formats. The MVP is specified with acceptance tests; every phase
has measurable entry and exit criteria, a risk register with early-warning
signals, and a research-to-engineering table that schedules hypotheses
H1--H7.

\begin{figure}[t]
\centering
% Roadmap timeline (weeks); width ~8.6cm; estimates. Phase 7 was added after Phases 0-6 were implemented.
\begin{tikzpicture}[x=0.197cm,y=0.5cm]
\foreach \w in {0,4,...,32}{\draw[black!15] (\w,0.3)--(\w,-7.7);\node[font=\tiny] at (\w,0.6){W\w};}
\foreach \ii/\lbl/\st/\en/\col in {
 0/{P0 Foundations}/0/3/navy,
 1/{P1 CPU MVP}/3/8/teal,
 2/{P2 Gate + loop}/8/13/amber,
 3/{P3 Fine-tuning}/12/18/violet,
 4/{P4 Eval + bench}/15/21/crimson,
 5/{P5 Hardening}/20/25/slate,
 6/{P6 Release}/25/28/leaf,
 7/{P7 CPU speed (added)}/28/32/teal}
 {\node[anchor=east,font=\tiny] at (-0.6,-\ii-0.2){\lbl};
  \fill[\col!75,rounded corners=1.5pt](\st,-\ii-0.5) rectangle (\en,-\ii+0.1);}
\foreach \x/\n in {3/G0,8/G1,13/G2,18/G3,21/G4,25/G5,28/G6}{\node[diamond,fill=black!75,inner sep=1.3pt] (g\n) at (\x,-8.1){};\node[font=\tiny,below=1pt] at (\x,-8.2){\n};}
\node[anchor=east,font=\tiny] at (-0.6,-8.1){gates};
\end{tikzpicture}

\caption{Roadmap timeline (weeks; estimates). Diamonds are decision gates
G0--G6 with go/no-go criteria defined in the roadmap document.}
\label{fig:gantt}
\end{figure}

\begin{figure}[t]
\centering
% Training (GPU) and serving (CPU) pipeline; width ~8.6cm
\begin{tikzpicture}[x=1cm,y=1cm]
\node[sblk=teal,text width=1.7cm]   (tr) at (0,0)   {Traces $+$ audits\\(verified only)};
\node[sblk=navy,text width=1.7cm]   (ds) at (2.2,0) {Dataset build\\filter, dedupe, 5\% abstain};
\node[sblk=violet,text width=1.7cm] (ft) at (4.4,0) {Unsloth QLoRA\\GPU (cloud or Colab)};
\node[sblk=amber,text width=1.7cm]  (gate) at (6.6,0) {Promotion gate\\held-out verifier};
\node[sblk=crimson,text width=1.7cm](exp) at (6.6,-1.5) {Merge $+$ export\\GGUF Q4\_K\_M};
\node[sblk=leaf,text width=1.7cm]   (srv) at (4.4,-1.5) {Serve on CPU\\llama-server or vLLM};
\node[sblk=slate,text width=1.7cm]  (hum) at (2.2,-1.5) {Human sign-off\\adapter hash};
\node[sblk=teal,text width=1.7cm]   (ev) at (0,-1.5) {\Harness\\regression suite};
\draw[arr] (tr)--(ds); \draw[arr] (ds)--(ft); \draw[arr] (ft)--(gate);
\draw[arr] (gate)--node[lab,right]{pass}(exp); \draw[arr] (exp)--(srv); \draw[arr] (srv)--(hum);
\draw[arr] (hum)--(ev);
\draw[dsh] (ev.south)--++(0,-0.45)-|node[lab,pos=.25,below]{new traces}(tr.south);
\draw[dsh] (gate.north)--++(0,0.45)-|node[lab,pos=.25,above]{fail: keep current adapter}(ft.north);
\end{tikzpicture}

\caption{Training and serving pipeline. Training runs on a rented GPU;
everything else runs on CPU.}
\label{fig:pipeline}
\end{figure}
